% =============================================================================
% 专利相似度计算过程详细文档
% Patent Similarity Calculation: Technical Documentation
% =============================================================================
\documentclass[12pt,a4paper]{article}

% ============================================================================
% 基础包
% ============================================================================
\usepackage[UTF8]{ctex}
\usepackage{geometry}
\usepackage{graphicx}
\usepackage{booktabs}
\usepackage{longtable}
\usepackage{amsmath,amssymb,amsfonts}
\usepackage{amsthm}
\usepackage{bm}
\usepackage{algorithm}
\usepackage{algpseudocode}
\usepackage{tikz}
\usepackage{float}
\usepackage{xcolor}
\usepackage{hyperref}
\usepackage{caption}
\usepackage{subcaption}
\usepackage{listings}
\usepackage{fancyhdr}
\usepackage{enumitem}
\usepackage{array}
\usepackage{multirow}
\usepackage{pgfplots}
\pgfplotsset{compat=1.18}
\usetikzlibrary{shapes,arrows,positioning,calc,fit,backgrounds,decorations.pathreplacing}

% ============================================================================
% 页面设置
% ============================================================================
\geometry{left=2.5cm,right=2.5cm,top=2.5cm,bottom=2.5cm}
\pagestyle{fancy}
\fancyhf{}
\fancyhead[L]{\leftmark}
\fancyhead[R]{\thepage}
\renewcommand{\headrulewidth}{0.4pt}

% ============================================================================
% 超链接设置
% ============================================================================
\hypersetup{
    colorlinks=true,
    linkcolor=blue,
    filecolor=magenta,
    urlcolor=cyan,
    citecolor=green,
    pdftitle={专利相似度计算过程详细文档},
    pdfauthor={专利相似度计算项目}
}

% ============================================================================
% 代码高亮设置
% ============================================================================
\definecolor{codegreen}{rgb}{0,0.6,0}
\definecolor{codegray}{rgb}{0.5,0.5,0.5}
\definecolor{codepurple}{rgb}{0.58,0,0.82}
\definecolor{backcolour}{rgb}{0.95,0.95,0.92}

\lstdefinestyle{mystyle}{
    backgroundcolor=\color{backcolour},
    commentstyle=\color{codegreen},
    keywordstyle=\color{magenta},
    numberstyle=\tiny\color{codegray},
    stringstyle=\color{codepurple},
    basicstyle=\ttfamily\footnotesize,
    breakatwhitespace=false,
    breaklines=true,
    captionpos=b,
    keepspaces=true,
    numbers=left,
    numbersep=5pt,
    showspaces=false,
    showstringspaces=false,
    showtabs=false,
    tabsize=2
}
\lstset{style=mystyle}

% ============================================================================
% 定理环境定义
% ============================================================================
\theoremstyle{definition}
\newtheorem{definition}{定义}[section]
\newtheorem{theorem}{定理}[section]
\newtheorem{lemma}{引理}[section]
\newtheorem{proposition}{命题}[section]
\newtheorem{example}{示例}[section]
\newtheorem{remark}{注}[section]
\newtheorem{algorithm_def}{算法}[section]

% ============================================================================
% 自定义命令
% ============================================================================
\newcommand{\vect}[1]{\mathbf{#1}}
\newcommand{\mat}[1]{\mathbf{#1}}
\newcommand{\norm}[1]{\left\|#1\right\|}
\newcommand{\inner}[2]{\left\langle #1, #2 \right\rangle}
\newcommand{\E}{\mathbb{E}}
\newcommand{\R}{\mathbb{R}}
\DeclareMathOperator{\cosine}{cosine}
\DeclareMathOperator{\softmax}{softmax}

% ============================================================================
% 文档信息
% ============================================================================
\title{\textbf{专利相似度计算过程详细文档}\\[0.3cm]
\large 基于 Sentence-BERT 的企业技术轨迹分析方法}
\author{专利相似度计算项目组}
\date{\today}

% ============================================================================
% 正文开始
% ============================================================================
\begin{document}

\maketitle
\thispagestyle{empty}

\begin{abstract}
本文档详细描述了基于 Sentence-BERT (SBERT) 的专利相似度计算完整流程。该流程包括数据预处理、语义嵌入生成、企业-年度聚合、以及多维度相似度计算等核心步骤。通过双模型交叉验证和引用加权稳健性检验，该方法能够有效识别企业技术转型信号，为投资研究、产业政策监测和竞争分析提供量化支撑。

\textbf{关键词：} 专利相似度；Sentence-BERT；语义嵌入；技术转型；余弦相似度
\end{abstract}

\tableofcontents
\newpage

% =============================================================================
\section{概述}
% =============================================================================

\subsection{研究目标}

本项目的核心目标是\textbf{量化企业每年专利组合与历史专利的语义相似度}，从而识别企业技术转型事件。具体而言，我们需要解决以下问题：

\begin{enumerate}
    \item 如何将专利文本（标题+摘要）转换为数值向量表示？
    \item 如何将同一企业同一年度的多个专利聚合成统一的年度技术表征？
    \item 如何计算当前年度与历史年度的技术相似度？
    \item 如何识别和验证技术转型事件？
\end{enumerate}

\subsection{计算流程概览}

整个计算流程分为四个主要阶段，如图~\ref{fig:pipeline}所示：

\begin{figure}[H]
    \centering
    \begin{tikzpicture}[
        node distance=2.5cm,
        stage/.style={rectangle, draw=blue!60, fill=blue!5, thick, minimum width=3cm, minimum height=1.2cm, align=center, rounded corners=3pt, font=\small},
        data/.style={rectangle, draw=green!60, fill=green!5, thick, minimum width=3cm, minimum height=1.2cm, align=center, rounded corners=3pt, font=\small},
        arrow/.style={->, >=stealth, thick, color=gray!70},
        label/.style={font=\footnotesize\itshape, color=gray!70}
    ]
    
    % 阶段节点
    \node[stage] (s1) {阶段1\\数据预处理};
    \node[stage, right of=s1] (s2) {阶段2\\SBERT编码};
    \node[stage, right of=s2] (s3) {阶段3\\Firm-Year聚合};
    \node[stage, right of=s3] (s4) {阶段4\\相似度计算};
    
    % 数据节点
    \node[data, below of=s1, node distance=2cm] (d1) {\texttt{patents.dta}\\原始专利数据};
    \node[data, below of=s2, node distance=2cm] (d2) {专利级嵌入\\$\vect{e}_i \in \R^{d}$};
    \node[data, below of=s3, node distance=2cm] (d3) {企业年度嵌入\\$\bar{\vect{e}}_{f,y} \in \R^{d}$};
    \node[data, below of=s4, node distance=2cm] (d4) {相似度指标\\$\cos(\bar{\vect{e}}_t, \bar{\vect{e}}_{t-k})$};
    
    % 箭头
    \draw[arrow] (d1) -- (s1);
    \draw[arrow] (s1) -- (s2) node[midway,above,label] {清洗文本};
    \draw[arrow] (s2) -- (s3) node[midway,above,label] {语义向量};
    \draw[arrow] (s3) -- (s4) node[midway,above,label] {年度表征};
    \draw[arrow] (s2) -- (d2);
    \draw[arrow] (s3) -- (d3);
    \draw[arrow] (s4) -- (d4);
    \draw[arrow] (d2) -- (s3);
    \draw[arrow] (d3) -- (s4);
    
    \end{tikzpicture}
    \caption{专利相似度计算流程概览}
    \label{fig:pipeline}
\end{figure}

\subsection{关键技术特点}

\begin{itemize}[leftmargin=2em]
    \item \textbf{双模型验证：}同时使用 MiniLM (384维) 和 DistilUSE (512维) 两种 SBERT 模型进行交叉验证，提高结果的可靠性。
    
    \item \textbf{双重聚合策略：}提供简单平均和引用加权平均两种聚合方式，后者给予高引用专利更高权重。
    
    \item \textbf{长文本处理：}针对超长专利文本，采用智能分块机制，优先按句子边界分割，必要时按 Token 边界分割。
    
    \item \textbf{多维度相似度：}计算与上一年、前三年均值、以及历史累积均值的多维度相似度指标。
\end{itemize}

% =============================================================================
\section{数据预处理}
% =============================================================================

\subsection{输入数据结构}

原始专利数据以 Stata 格式 (\texttt{.dta}) 存储，包含以下核心字段：

\begin{table}[H]
    \centering
    \caption{输入数据字段说明}
    \label{tab:input-fields}
    \begin{tabular}{@{}llcp{6cm}@{}}
        \toprule
        \textbf{字段名} & \textbf{类型} & \textbf{必需} & \textbf{说明} \\
        \midrule
        \texttt{stkcd} & 字符串 & 是 & 公司股票代码，企业唯一标识 \\
        \texttt{p\_year} & 整数 & 是 & 专利申请年份 \\
        \texttt{p\_tt} & 字符串 & 是 & 专利标题 \\
        \texttt{p\_abs} & 字符串 & 是 & 专利摘要 \\
        \texttt{p\_id} & 字符串 & 否 & 专利唯一标识符 \\
        \texttt{p\_cite} & 数值 & 否 & 被引证次数，用于加权聚合 \\
        \texttt{p\_date} & 日期 & 否 & 专利申请日期 \\
        \texttt{p\_type} & 字符串 & 否 & 专利类型（发明/实用新型等） \\
        \texttt{p\_ipc} & 字符串 & 否 & IPC 分类号 \\
        \bottomrule
    \end{tabular}
\end{table}

\subsection{预处理步骤}

\begin{algorithm}[H]
    \caption{数据预处理流程}
    \label{alg:preprocessing}
    \begin{algorithmic}[1]
        \Require 原始专利数据表 $D$
        \Ensure 预处理后的数据表 $D'$
        \State 验证必需字段存在性：\texttt{stkcd}, \texttt{p\_year}, \texttt{p\_tt}, \texttt{p\_abs}
        \State 处理缺失值：\texttt{p\_cite} $\leftarrow$ \texttt{fillna(0)}
        \State 合并文本字段：
        \State \quad $D'.\text{\texttt{combined\_text}} \leftarrow D.\text{\texttt{p\_tt}} + \text{`` ''} + D.\text{\texttt{p\_abs}}$
        \State 标记空文本：
        \State \quad $D'.\text{\texttt{has\_text}} \leftarrow \mathbb{I}(D'.\text{\texttt{combined\_text}}.\text{str.len()} > 0)$
        \State 创建复合键：
        \State \quad $D'.\text{\texttt{stkcd\_year}} \leftarrow D.\text{\texttt{stkcd}} + \text{``\_''} + D.\text{\texttt{p\_year}}.\text{astype(str)}$
        \State \Return $D'$
    \end{algorithmic}
\end{algorithm}

\subsection{示例数据}

以公司2（股票代码：2）2002年的专利数据为例：

\begin{example}[数据预处理示例]
\label{ex:preprocessing}
考虑公司2在2002年申请的一项专利：

\begin{itemize}
    \item \textbf{股票代码 (stkcd):} 2
    \item \textbf{年份 (p\_year):} 2002
    \item \textbf{标题 (p\_tt):} 情景花园洋房
    \item \textbf{摘要 (p\_abs):} 一种情景花园洋房，每单元两户，每户不仅包含了室内空间，还包含了与外界情景环境容为一体的室外空间...
    \item \textbf{引用次数 (p\_cite):} 12
\end{itemize}

预处理后：
\begin{itemize}
    \item \textbf{combined\_text:} ``情景花园洋房 一种情景花园洋房，每单元两户...''
    \item \textbf{has\_text:} \texttt{True}
    \item \textbf{stkcd\_year:} ``2\_2002''
\end{itemize}
\end{example}

% =============================================================================
\section{语义嵌入生成}
% =============================================================================

\subsection{Sentence-BERT 模型}

Sentence-BERT (SBERT) \cite{reimers2019sentence} 是一种基于 BERT 的句子嵌入模型，通过对 Siamese 网络结构进行微调，能够生成语义丰富的句子向量。与传统 BERT 相比，SBERT 的优势在于：

\begin{enumerate}
    \item \textbf{计算效率高：}句子对的相似度计算从 $O(n^2)$ 降低到 $O(n)$
    \item \textbf{语义捕获能力强：}通过对比学习，向量空间具有更好的语义结构
    \item \textbf{多语言支持：}预训练模型支持包括中文在内的50+种语言
\end{enumerate}

\subsection{使用的模型}

本项目使用以下两种 SBERT 模型进行交叉验证：

\begin{table}[H]
    \centering
    \caption{SBERT 模型规格对比}
    \label{tab:models}
    \begin{tabular}{@{}lcc@{}}
        \toprule
        \textbf{特性} & \textbf{MiniLM} & \textbf{DistilUSE} \\
        \midrule
        模型名称 & paraphrase-multilingual-MiniLM-L12-v2 & distiluse-base-multilingual-cased-v2 \\
        基础架构 & MiniLM-L12 (蒸馏BERT) & DistilBERT + 全连接层 \\
        输出维度 & 384 & 512 \\
        最大序列长度 & 512 tokens & 512 tokens \\
        支持语言 & 50+ & 50+ \\
        模型简称 & \texttt{minilm} & \texttt{distiluse} \\
        \bottomrule
    \end{tabular}
\end{table}

\subsection{嵌入计算过程}

\begin{definition}[专利嵌入]
给定专利 $i$ 的合并文本 $t_i$，SBERT 模型 $\mathcal{M}$ 将其映射为 $d$ 维向量：
\begin{equation}
    \vect{e}_i = \mathcal{M}(t_i) \in \R^d
\end{equation}
其中 $d = 384$ (MiniLM) 或 $d = 512$ (DistilUSE)。
\end{definition}

嵌入计算的具体过程如下：

\begin{algorithm}[H]
    \caption{SBERT 嵌入计算}
    \label{alg:embedding}
    \begin{algorithmic}[1]
        \Require 专利文本列表 $T = [t_1, t_2, ..., t_n]$，SBERT模型 $\mathcal{M}$，批次大小 $B$
        \Ensure 嵌入矩阵 $\mat{E} \in \R^{n \times d}$
        \State 初始化空列表 $E \leftarrow []$
        \For{每个批次 $b$ in $\text{batch}(T, B)$}
            \State $\mat{E}_b \leftarrow \mathcal{M}.\text{encode}(b)$
            \State $E.\text{append}(\mat{E}_b)$
        \EndFor
        \State $\mat{E} \leftarrow \text{concatenate}(E)$
        \State \Return $\mat{E}$
    \end{algorithmic}
\end{algorithm}

\subsection{长文本处理机制}

由于部分专利文本（特别是包含详细技术描述的摘要）可能超过模型的最大序列长度（512 tokens），需要特殊的处理机制：

\begin{algorithm}[H]
    \caption{长文本分块处理}
    \label{alg:chunking}
    \begin{algorithmic}[1]
        \Require 文本 $t$，最大长度 $L_{\max} = 512$
        \Ensure 嵌入向量 $\vect{e}$
        \State 对 $t$ 进行 Tokenization，得到 Token 序列 $[w_1, w_2, ..., w_m]$
        \If{$m \leq L_{\max}$}
            \State \Return $\mathcal{M}(t)$ \Comment{标准编码}
        \Else
            \State 按句子边界分割：$S \leftarrow \text{split\_by\_sentences}(t)$
            \State 初始化空列表 $C \leftarrow []$
            \State 初始化当前块 $c \leftarrow \text{``''}$
            \For{每个句子 $s$ in $S$}
                \If{$\text{len}(c) + \text{len}(s) \leq L_{\max}$}
                    \State $c \leftarrow c + s$
                \Else
                    \State $C.\text{append}(c)$
                    \State $c \leftarrow s$
                \EndIf
            \EndFor
            \State $C.\text{append}(c)$ \Comment{添加最后一个块}
            \State \Return $\frac{1}{|C|} \sum_{c \in C} \mathcal{M}(c)$ \Comment{块嵌入均值聚合}
        \EndIf
    \end{algorithmic}
\end{algorithm}

句子分割使用正则表达式：
\begin{equation}
    \text{pattern} = (?<=[\text{。；}!?！？;])\s+|\n+
\end{equation}

% =============================================================================
\section{企业-年度聚合}
% =============================================================================

\subsection{聚合目标}

将同一企业（\texttt{stkcd}）同一年度（\texttt{p\_year}）的所有专利嵌入聚合成单一的年度技术表征向量。这涉及两个核心问题：

\begin{enumerate}
    \item \textbf{聚合方法：}如何组合多个专利向量？
    \item \textbf{空文本处理：}是否将无文本专利纳入聚合？
\end{enumerate}

\subsection{简单平均聚合}

\begin{definition}[简单平均嵌入]
对于企业 $f$ 在年份 $y$ 的专利集合 $\mathcal{P}_{f,y} = \{1, 2, ..., N_{f,y}\}$，简单平均嵌入定义为：
\begin{equation}
    \bar{\vect{e}}_{f,y} = \frac{1}{N_{f,y}} \sum_{i \in \mathcal{P}_{f,y}} \vect{e}_i
\end{equation}
其中 $N_{f,y}$ 是专利总数（包括空文本专利，其嵌入为零向量）。
\end{definition}

\subsection{引用加权聚合}

引用次数反映专利的技术影响力，高引用专利可能更好地代表企业的核心技术方向。

\begin{definition}[引用加权嵌入]
设专利 $i$ 的引用次数为 $c_i$，则引用加权嵌入为：
\begin{equation}
    \tilde{\vect{e}}_{f,y} = \frac{1}{\sum_{i \in \mathcal{P}_{f,y}} c_i} \sum_{i \in \mathcal{P}_{f,y}} c_i \cdot \vect{e}_i
\end{equation}
\end{definition}

\subsection{聚合算法}

\begin{algorithm}[H]
    \caption{企业-年度聚合}
    \label{alg:aggregation}
    \begin{algorithmic}[1]
        \Require 专利嵌入矩阵 $\mat{E}$，元数据表 $M$（含 \texttt{stkcd}, \texttt{p\_year}, \texttt{p\_cite}, \texttt{has\_text}）
        \Require 是否包含空文本标志 $\text{include\_empty}$
        \Ensure 企业-年度嵌入表 $A$
        \State 按 \texttt{stkcd\_year} 分组：$G \leftarrow M.\text{groupby}(\text{`stkcd\_year'})$
        \For{每个组 $g$ in $G$}
            \State 提取企业代码：$f \leftarrow g.\text{stkcd}.\text{iloc}[0]$
            \State 提取年份：$y \leftarrow g.\text{p\_year}.\text{iloc}[0]$
            \State 获取专利索引：$I \leftarrow g.\text{index}$
            \If{not $\text{include\_empty}$}
                \State $I \leftarrow I[g.\text{has\_text} == \text{True}]$
            \EndIf
            \State 提取嵌入：$\mat{E}_g \leftarrow \mat{E}[I, :]$
            \State 计算简单平均：$\bar{\vect{e}} \leftarrow \text{mean}(\mat{E}_g, \text{axis}=0)$
            \State 计算引用加权：$\tilde{\vect{e}} \leftarrow \text{citation\_weighted\_mean}(\mat{E}_g, g.\text{p\_cite})$
            \State 记录元数据：
            \State \quad $n_{\text{patents}} \leftarrow \text{len}(g)$
            \State \quad $n_{\text{texts}} \leftarrow \text{sum}(g.\text{has\_text})$
            \State \quad $c_{\text{total}} \leftarrow \text{sum}(g.\text{p\_cite})$
            \State 存储结果到 $A$
        \EndFor
        \State \Return $A$
    \end{algorithmic}
\end{algorithm}

\subsection{聚合示例}

\begin{example}[公司2，2002年聚合]
假设公司2在2002年有3项专利，其嵌入和引用数据如下：

\begin{table}[H]
    \centering
    \begin{tabular}{@{}cccc@{}}
        \toprule
        \textbf{专利ID} & \textbf{嵌入} $\vect{e}_i$ & \textbf{引用} $c_i$ & \textbf{是否空文本} \\
        \midrule
        P1 & $[0.1, -0.2, ..., 0.05]$ & 12 & 否 \\
        P2 & $[0.05, 0.1, ..., -0.1]$ & 7 & 否 \\
        P3 & $[-0.05, 0.15, ..., 0.02]$ & 9 & 否 \\
        \bottomrule
    \end{tabular}
\end{table}

则聚合结果为：
\begin{align}
    \bar{\vect{e}} &= \frac{\vect{e}_1 + \vect{e}_2 + \vect{e}_3}{3} \\
    \tilde{\vect{e}} &= \frac{12\vect{e}_1 + 7\vect{e}_2 + 9\vect{e}_3}{12 + 7 + 9} = \frac{12\vect{e}_1 + 7\vect{e}_2 + 9\vect{e}_3}{28}
\end{align}

元数据：
\begin{itemize}
    \item $n_{\text{patents}} = 3$
    \item $n_{\text{texts}} = 3$
    \item $c_{\text{total}} = 28$
    \item $c_{\text{mean}} = 28/3 \approx 9.33$
\end{itemize}
\end{example}

% =============================================================================
\section{相似度计算}
% =============================================================================

\subsection{余弦相似度定义}

\begin{definition}[余弦相似度]
对于两个非零向量 $\vect{u}, \vect{v} \in \R^d$，其余弦相似度定义为：
\begin{equation}
    \cosine(\vect{u}, \vect{v}) = \frac{\inner{\vect{u}}{\vect{v}}}{\norm{\vect{u}} \cdot \norm{\vect{v}}} = \frac{\sum_{i=1}^{d} u_i v_i}{\sqrt{\sum_{i=1}^{d} u_i^2} \cdot \sqrt{\sum_{i=1}^{d} v_i^2}}
\end{equation}
余弦相似度的取值范围为 $[-1, 1]$，其中：
\begin{itemize}
    \item $1$ 表示方向完全相同（最相似）
    \item $0$ 表示正交（无相关性）
    \item $-1$ 表示方向完全相反
\end{itemize}
\end{definition}

\begin{proposition}[归一化向量的余弦相似度]
如果 $\vect{u}$ 和 $\vect{v}$ 都是单位向量（即 $\norm{\vect{u}} = \norm{\vect{v}} = 1$），则：
\begin{equation}
    \cosine(\vect{u}, \vect{v}) = \inner{\vect{u}}{\vect{v}}
\end{equation}
\end{proposition}

\subsection{三种相似度指标}

为了从不同时间维度衡量企业技术方向的连续性，我们计算以下三种相似度指标：

\subsubsection{滞后一年相似度 (cos\_sim\_lag1)}

衡量当前年度与上一年度技术方向的偏离程度：
\begin{equation}
    \text{cos\_sim\_lag1}_{f,y} = \cosine(\bar{\vect{e}}_{f,y}, \bar{\vect{e}}_{f,y-1})
\end{equation}

\subsubsection{滞后三年相似度 (cos\_sim\_lag3)}

衡量当前年度与前三年平均技术方向的偏离程度：
\begin{equation}
    \text{cos\_sim\_lag3}_{f,y} = \cosine\left(\bar{\vect{e}}_{f,y}, \frac{1}{3}\sum_{k=1}^{3} \bar{\vect{e}}_{f,y-k}\right)
\end{equation}
当企业在前三年内任一年份无专利时，该指标为 \texttt{NA}。

\subsubsection{累积相似度 (cos\_sim\_cumulative)}

衡量当前年度与历史累积技术轨迹的偏离程度：
\begin{equation}
    \text{cos\_sim\_cumulative}_{f,y} = \cosine\left(\bar{\vect{e}}_{f,y}, \frac{1}{y - y_{\min}}\sum_{k=y_{\min}}^{y-1} \bar{\vect{e}}_{f,k}\right)
\end{equation}
其中 $y_{\min}$ 是企业最早有专利记录的年份。

\subsection{相似度计算算法}

\begin{algorithm}[H]
    \caption{计算企业年度相似度}
    \label{alg:similarity}
    \begin{algorithmic}[1]
        \Require 企业-年度嵌入表 $A$（已按 \texttt{stkcd} 和 \texttt{p\_year} 排序）
        \Ensure 相似度结果表 $S$
        \State 初始化空列表 $S$
        \For{每个企业 $f$ in $A.\text{stkcd}.\text{unique}()$}
            \State 提取企业数据：$A_f \leftarrow A[A.\text{stkcd} == f]$
            \State 按年份排序：$A_f \leftarrow A_f.\text{sort\_values}(\text{`p\_year'})$
            \State 获取年份列表：$Y \leftarrow A_f.\text{p\_year}.\text{values}$
            \State 获取嵌入矩阵：$\mat{E}_f \leftarrow A_f[\text{emb\_cols}].\text{values}$
            \For{$i = 0$ to $\text{len}(Y) - 1$}
                \State $y \leftarrow Y[i]$
                \State $\vect{e}_t \leftarrow \mat{E}_f[i, :]$ \Comment{当前年度嵌入}
                \State 初始化行记录 $row \leftarrow \{\text{stkcd}: f, \text{p\_year}: y\}$
                
                \State \Comment{计算滞后一年相似度}
                \If{$i > 0$}
                    \State $\vect{e}_{t-1} \leftarrow \mat{E}_f[i-1, :]$
                    \State $row[\text{cos\_sim\_lag1}] \leftarrow \cosine(\vect{e}_t, \vect{e}_{t-1})$
                \Else
                    \State $row[\text{cos\_sim\_lag1}] \leftarrow \text{NA}$
                \EndIf
                
                \State \Comment{计算滞后三年相似度}
                \If{$i \geq 3$}
                    \State $\bar{\vect{e}}_{t-3:t-1} \leftarrow \text{mean}(\mat{E}_f[i-3:i, :], \text{axis}=0)$
                    \State $row[\text{cos\_sim\_lag3}] \leftarrow \cosine(\vect{e}_t, \bar{\vect{e}}_{t-3:t-1})$
                \Else
                    \State $row[\text{cos\_sim\_lag3}] \leftarrow \text{NA}$
                \EndIf
                
                \State \Comment{计算累积相似度}
                \If{$i > 0$}
                    \State $\bar{\vect{e}}_{\text{cum}} \leftarrow \text{mean}(\mat{E}_f[0:i, :], \text{axis}=0)$
                    \State $row[\text{cos\_sim\_cumulative}] \leftarrow \cosine(\vect{e}_t, \bar{\vect{e}}_{\text{cum}})$
                \Else
                    \State $row[\text{cos\_sim\_cumulative}] \leftarrow \text{NA}$
                \EndIf
                
                \State $S.\text{append}(row)$
            \EndFor
        \EndFor
        \State \Return $\text{DataFrame}(S)$
    \end{algorithmic}
\end{algorithm}

\subsection{安全计算机制}

在实际计算中需要处理以下边界情况：

\begin{itemize}
    \item \textbf{零向量检测：}如果任一向量为零向量（全零），返回 \texttt{NA}
    \item \textbf{NaN/Inf 检测：}如果向量包含 NaN 或 Inf 值，返回 \texttt{NA}
    \item \textbf{分组隔离：}确保不同企业之间的数据不交叉计算
\end{itemize}

% =============================================================================
\section{技术转型识别}
% =============================================================================

\subsection{转型识别标准}

基于相似度指标，我们定义以下技术转型识别标准：

\begin{definition}[技术转型事件]
企业 $f$ 在年份 $y$ 发生技术转型，当满足：
\begin{equation}
    \text{cos\_sim\_lag1}_{f,y} < \theta \quad \text{且} \quad n_{\text{patents},f,y} \geq n_{\min}
\end{equation}
其中 $\theta = 0.5$ 为相似度阈值，$n_{\min} = 5$ 为最小专利数要求。
\end{definition}

双模型交叉验证要求：
\begin{equation}
    \text{cos\_sim\_lag1}_{f,y}^{\text{MiniLM}} < \theta \quad \text{且} \quad \text{cos\_sim\_lag1}_{f,y}^{\text{DistilUSE}} < \theta
\end{equation}

\subsection{转型模式分类}

根据相似度轨迹特征，技术转型可分为以下模式：

\begin{table}[H]
    \centering
    \caption{技术转型模式分类}
    \label{tab:transformation-patterns}
    \begin{tabular}{@{}lp{5cm}p{5cm}@{}}
        \toprule
        \textbf{转型模式} & \textbf{特征描述} & \textbf{典型案例} \\
        \midrule
        渐进式转型 & 相似度逐年下降，经历2--3年转型期 & 公司2 (2007--2010)：从建筑设计转向电子控制 \\
        战略式转型 & 长期高度稳定后突然断裂 & 公司12 (2020--2021)：从电子玻璃转向光伏 \\
        探索式转型 & 快速恢复并达到更高相似度 & 公司518：转型后技术方向更聚焦 \\
        震荡式转型 & 相似度波动，多次转型信号 & 某些企业可能经历多次方向调整 \\
        \bottomrule
    \end{tabular}
\end{table}

\subsection{案例验证}

\subsubsection{案例1：公司2（渐进式转型）}

公司2（股票代码：2）是房地产行业相关企业，数据跨度 2002--2014 年。

\begin{table}[H]
    \centering
    \caption{公司2相似度轨迹（部分年份）}
    \label{tab:company2}
    \begin{tabular}{@{}ccccc@{}}
        \toprule
        \textbf{年份} & \textbf{专利数} & \textbf{MiniLM} & \textbf{DistilUSE} & \textbf{转型信号} \\
        \midrule
        2002 & 3 & -- & -- & -- \\
        2004 & 9 & 0.723 & 0.689 & -- \\
        2005 & 2 & 0.891 & 0.856 & -- \\
        2006 & 41 & 0.817 & 0.799 & -- \\
        2007 & 3 & 0.626 & 0.487 & -- \\
        2008 & 30 & 0.606 & 0.409 & -- \\
        \rowcolor{red!15}
        2009 & 1 & 0.334 & 0.182 & \checkmark \\
        \rowcolor{red!15}
        2010 & 12 & 0.378 & 0.233 & \checkmark \\
        2013 & 2 & 0.732 & 0.666 & -- \\
        \bottomrule
    \end{tabular}
\end{table}

\textbf{转型特征分析：}
\begin{enumerate}
    \item \textbf{渐进下降：}2007--2009 年相似度从 0.8 逐步下降至 0.3
    \item \textbf{转型断裂：}2009--2010 年相似度低于 0.5，触发转型信号
    \item \textbf{恢复稳定：}2013 年后相似度回升至 0.7+
\end{enumerate}

\textbf{专利文本对比：}

\begin{itemize}
    \item \textbf{转型前（2006年）：}``带垃圾收集装置的橱柜''——传统家居设计
    \item \textbf{转型期（2009年）：}``止吠器''——电子控制/微电路技术
\end{itemize}

从传统建筑设计转向电子控制技术，技术方向发生根本性变化。

\subsubsection{案例2：公司12（战略式转型）}

公司12（股票代码：12）是玻璃制造行业企业，数据跨度 2001--2023 年。

\begin{table}[H]
    \centering
    \caption{公司12相似度轨迹（转型前后）}
    \label{tab:company12}
    \begin{tabular}{@{}ccccc@{}}
        \toprule
        \textbf{年份} & \textbf{专利数} & \textbf{MiniLM} & \textbf{DistilUSE} & \textbf{转型信号} \\
        \midrule
        2017 & 70 & 0.859 & 0.881 & -- \\
        2018 & 47 & 0.940 & 0.927 & -- \\
        2019 & 12 & 0.954 & 0.914 & -- \\
        \rowcolor{red!15}
        2020 & 1 & 0.493 & 0.122 & \checkmark \\
        \rowcolor{red!15}
        2021 & 7 & 0.406 & 0.131 & \checkmark \\
        2022 & 3 & 0.820 & 0.673 & -- \\
        \bottomrule
    \end{tabular}
\end{table}

\textbf{转型特征分析：}
\begin{enumerate}
    \item \textbf{长期深耕：}2007--2019 年相似度稳定在 0.7--0.95 之间
    \item \textbf{突然断裂：}2020 年从 0.95 骤降至 0.49/0.12
    \item \textbf{战略转型：}2021 年进一步降至 0.41/0.13
    \item \textbf{快速稳定：}2022 年后相似度回升，新方向趋于稳定
\end{enumerate}

\textbf{专利文本对比：}
\begin{itemize}
    \item \textbf{转型前（2019年）：}``电容式触摸屏''——电子玻璃/触控屏
    \item \textbf{转型期（2021年）：}``光伏组件测试架''——光伏/新能源
\end{itemize}

典型的``战略式转型''：在电子玻璃领域深耕多年后，主动切换至光伏新能源赛道。

% =============================================================================
\section{实现细节}
% =============================================================================

\subsection{计算环境}

\begin{table}[H]
    \centering
    \caption{计算环境配置}
    \label{tab:environment}
    \begin{tabular}{@{}ll@{}}
        \toprule
        \textbf{组件} & \textbf{配置/版本} \\
        \midrule
        编程语言 & Python 3.8+ / R 4.0+ \\
        深度学习框架 & PyTorch $\geq$ 2.0.0 \\
        Transformer库 & Transformers $\geq$ 4.30.0 \\
        SBERT库 & sentence-transformers $\geq$ 2.2.0 \\
        数据处理 & pandas $\geq$ 2.0.0, numpy $\geq$ 1.24.0 \\
        计算设备 & GPU (CUDA) 推荐，支持多GPU \\
        \bottomrule
    \end{tabular}
\end{table}

\subsection{关键代码片段}

\subsubsection{数据预处理 (Python)}

\begin{lstlisting}[language=Python, caption=数据预处理核心代码]
def load_and_prepare_data(file_path: str) -> pd.DataFrame:
    """加载并预处理专利数据"""
    df = pd.read_stata(file_path)
    
    # 合并标题和摘要
    df['combined_text'] = df['p_tt'].fillna('') + ' ' + df['p_abs'].fillna('')
    df['has_text'] = df['combined_text'].str.len() > 0
    
    # 创建复合键
    df['stkcd_year'] = df['stkcd'].astype(str) + '_' + df['p_year'].astype(str)
    
    # 填充缺失引用
    df['p_cite'] = df['p_cite'].fillna(0)
    
    return df
\end{lstlisting}

\subsubsection{SBERT编码 (Python)}

\begin{lstlisting}[language=Python, caption=SBERT编码核心代码]
from sentence_transformers import SentenceTransformer

class SBertEmbedder:
    def __init__(self, model_path: str, device: str = 'cuda'):
        self.model = SentenceTransformer(model_path, device=device)
        self.max_seq_length = self.model.max_seq_length
    
    def encode(self, texts: List[str], batch_size: int = 256) -> np.ndarray:
        """批量编码文本"""
        embeddings = self.model.encode(
            texts,
            batch_size=batch_size,
            show_progress_bar=True,
            convert_to_numpy=True
        )
        return embeddings
\end{lstlisting}

\subsubsection{相似度计算 (R)}

\begin{lstlisting}[language=R, caption=R语言相似度计算核心代码]
calculate_similarities <- function(embeddings_df) {
  # 按 stkcd 分组计算
  results <- embeddings_df %>%
    group_by(stkcd) %>%
    arrange(p_year) %>%
    mutate(
      # 滞后一年相似度
      cos_sim_lag1 = calculate_cos_sim(emb_matrix, lag = 1),
      
      # 滞后三年相似度
      cos_sim_lag3 = calculate_cos_sim(emb_matrix, lag = 3, average = TRUE),
      
      # 累积相似度
      cos_sim_cumulative = calculate_cumulative_cos_sim(emb_matrix)
    ) %>%
    ungroup()
  
  return(results)
}
\end{lstlisting}

\subsection{性能优化}

\begin{enumerate}
    \item \textbf{GPU加速：}使用 CUDA 加速 SBERT 编码，比 CPU 快 10-50 倍
    \item \textbf{批次处理：}默认批次大小 256，可根据 GPU 内存调整
    \item \textbf{多GPU支持：}支持多 GPU 并行处理，通过 \texttt{--multi-gpu} 启用
    \item \textbf{内存管理：}流式处理大数据集，避免一次性加载全部数据
\end{enumerate}

% =============================================================================
\section{输出数据结构}
% =============================================================================

\subsection{企业-年度嵌入文件}

文件名格式：\texttt{stkcd\_year\_\{model\}\_embeddings.csv}

\begin{table}[H]
    \centering
    \caption{企业-年度嵌入文件字段说明}
    \label{tab:output-embeddings}
    \begin{tabular}{@{}llp{7cm}@{}}
        \toprule
        \textbf{字段名} & \textbf{类型} & \textbf{说明} \\
        \midrule
        \texttt{stkcd} & 字符串 & 企业股票代码 \\
        \texttt{p\_year} & 整数 & 年份 \\
        \texttt{stkcd\_year} & 字符串 & 复合键 (stkcd\_p\_year) \\
        \texttt{n\_patents} & 整数 & 该年度专利总数 \\
        \texttt{n\_texts\_used} & 整数 & 用于聚合的非空文本专利数 \\
        \texttt{total\_citations} & 整数 & 引用总数 \\
        \texttt{mean\_citations} & 浮点数 & 平均引用次数 \\
        \texttt{emb\_0} -- \texttt{emb\_N} & 浮点数 & 嵌入向量分量 (N=383 或 511) \\
        \bottomrule
    \end{tabular}
\end{table}

\subsection{相似度结果文件}

文件名格式：\texttt{stkcd\_year\_similarity\_\{model\}.csv}

\begin{table}[H]
    \centering
    \caption{相似度结果文件字段说明}
    \label{tab:output-similarity}
    \begin{tabular}{@{}llp{7cm}@{}}
        \toprule
        \textbf{字段名} & \textbf{类型} & \textbf{说明} \\
        \midrule
        \texttt{stkcd} & 字符串 & 企业股票代码 \\
        \texttt{p\_year} & 整数 & 年份 \\
        \texttt{n\_patents} & 整数 & 该年度专利数 \\
        \texttt{cos\_sim\_lag1} & 浮点数 & 与上一年相似度 \\
        \texttt{cos\_sim\_lag3} & 浮点数 & 与前三年均值相似度 \\
        \texttt{cos\_sim\_cumulative} & 浮点数 & 与历史累积均值相似度 \\
        \bottomrule
    \end{tabular}
\end{table}

% =============================================================================
\section{质量控制与验证}
% =============================================================================

\subsection{双模型交叉验证}

为确保结果的可靠性，我们使用两个独立的 SBERT 模型进行交叉验证：

\begin{equation}
    \text{一致性得分} = \frac{1}{N} \sum_{i=1}^{N} \mathbb{I}\left(|\text{sim}_{i}^{\text{MiniLM}} - \text{sim}_{i}^{\text{Distil}}| < \epsilon\right)
\end{equation}

经验表明，两个模型计算的相似度高度一致（相关系数 $> 0.9$），异常差异通常指向数据质量问题。

\subsection{引用加权稳健性检验}

通过比较简单平均和引用加权的结果，检验相似度指标对聚合方法的敏感性：

\begin{equation}
    \Delta_{f,y} = \text{cos\_sim}_{f,y}^{\text{simple}} - \text{cos\_sim}_{f,y}^{\text{citation-weighted}}
\end{equation}

如果 $|\Delta_{f,y}|$ 过大，表明该年度专利质量分布不均，需要审慎解读。

\subsection{人工验证}

对识别出的转型案例进行人工专利文本审查，验证：
\begin{enumerate}
    \item 转型前后专利技术领域是否确实发生变化
    \item 相似度下降是否与技术方向转变一致
    \item 是否存在数据质量问题（如文本缺失、编码错误）
\end{enumerate}

% =============================================================================
\section{总结}
% =============================================================================

本文档详细描述了基于 Sentence-BERT 的专利相似度计算完整流程。该方法通过以下步骤实现企业技术轨迹的量化分析：

\begin{enumerate}
    \item \textbf{数据预处理：}清洗原始专利数据，合并标题和摘要，创建企业-年度复合键。
    
    \item \textbf{语义嵌入：}使用 SBERT 模型将专利文本转换为高维向量，采用智能分块机制处理超长文本。
    
    \item \textbf{企业-年度聚合：}提供简单平均和引用加权两种聚合策略，生成企业年度技术表征向量。
    
    \item \textbf{相似度计算：}计算滞后一年、滞后三年和累积三种相似度指标，多维度衡量技术连续性。
    
    \item \textbf{转型识别：}基于相似度阈值和双模型交叉验证识别技术转型事件，分类转型模式。
\end{enumerate}

该方法的优势在于：
\begin{itemize}
    \item \textbf{语义层面：}基于 SBERT 捕获专利文本深层语义，超越传统关键词匹配
    \item \textbf{技术完备：}长文本处理、多GPU加速、流式分块，支持大规模数据
    \item \textbf{结果可靠：}双模型交叉验证 + 引用加权稳健性检验
\end{itemize}

应用价值包括投资研究（识别潜在的技术转型标的）、产业政策监测（跟踪行业技术演进）和竞争分析（预警竞争对手技术方向变化）。

% =============================================================================
\begin{thebibliography}{99}
\bibitem{reimers2019sentence}
Reimers, N., \& Gurevych, I. (2019). Sentence-BERT: Sentence embeddings using Siamese BERT-networks. \textit{Proceedings of the 2019 Conference on Empirical Methods in Natural Language Processing}, 3982--3992.

\bibitem{devlin2018bert}
Devlin, J., Chang, M. W., Lee, K., \& Toutanova, K. (2018). BERT: Pre-training of deep bidirectional transformers for language understanding. \textit{Proceedings of NAACL-HLT}, 4171--4186.

\bibitem{vaswani2017attention}
Vaswani, A., Shazeer, N., Parmar, N., Uszkoreit, J., Jones, L., Gomez, A. N., ... \& Polosukhin, I. (2017). Attention is all you need. \textit{Advances in Neural Information Processing Systems}, 30.
\end{thebibliography}

\end{document}
